%! Dividing Polynomials; Remainder and Factor Theorems — Unit 4, Lesson 4.3 — Group Activity
\documentclass[10pt]{article}
\usepackage{algebra2-article}
\usepackage{algebra2-boxes}
\newcommand{\TallMath}[1]{$\displaystyle #1\rule[-1.4em]{0pt}{3.2em}$}

\begin{document}

\pageheader{Unit 4, Lesson 4.3}{Group Activity}
\namepartnerperiod

\vspace{0.08in}

\begin{headlinebox}{forestbg}
\small\textbf{Divide and Conquer.} Every problem today runs on one statement:
\textbf{dividend $=$ divisor $\times$ quotient $+$ remainder}. Three habits keep you out of trouble:
write the dividend in standard form with a \emph{placeholder $0$} for every missing power; remember
that $x + 2$ means $r = -2$; and read the remainder --- a remainder of $0$ is the Factor Theorem
telling you that you just found a factor.
\end{headlinebox}

\vspace{0.06in}

\begin{tcolorbox}[colback=white, colframe=black!40, title=\textbf{Tier R --- Remediate},
    fonttitle=\bfseries, arc=1.5mm, left=3mm, right=3mm, top=2mm, bottom=2mm, breakable]
\begin{enumerate}[leftmargin=1.7em, itemsep=9pt]
  \item \textbf{Monomial divisor.} Split the fraction term by term.

    \vspace{2pt}
    $\dfrac{10x^4 - 15x^3 + 5x^2}{5x^2} = $ \blank{5.0cm}

    \vspace{2pt}
    Careful with the last term --- it divides out to \blank{1.2cm}, not $0$.

  \item \textbf{Long division, step by step.} Divide $x^2 + 7x + 12$ by $x + 3$.

    \vspace{3pt}
    \emph{Step 1 --- first term of the quotient:} $\dfrac{x^2}{x} = $ \blank{1.4cm}

    \vspace{3pt}
    \emph{Step 2 --- multiply and subtract:} $x(x+3) = $ \blank{2.4cm}, and
    $(x^2+7x) - ($\blank{2.4cm}$) = $ \blank{1.6cm}

    \vspace{3pt}
    \emph{Step 3 --- bring down the $12$ and repeat:} $\dfrac{4x}{x} = $ \blank{1.0cm}, \;
    $4(x+3) = $ \blank{2.0cm}

    \vspace{3pt}
    \emph{Step 4:} Quotient \blank{2.0cm}, remainder \blank{1.2cm}. \; Check by multiplying:
    $(x+3)($\blank{1.6cm}$) = x^2+7x+12$? \blank{1.2cm}

  \item \textbf{Synthetic division.} Divide $x^3 + 3x^2 - 4x - 12$ by $x + 3$.

    \vspace{3pt}
    First, the sign flip: the divisor is $x + 3$, so $r = $ \blank{1.0cm}
    \begin{center}
    \renewcommand{\arraystretch}{1.6}
    \begin{tabular}{r|C{1.4cm}C{1.4cm}C{1.4cm}C{1.4cm}}
        & $1$ & $3$ & $-4$ & $-12$ \\
        &     &     &      &       \\ \cline{2-5}
        &     &     &      &       \\
    \end{tabular}
    \end{center}

    \vspace{2pt}
    Quotient: \blank{3.0cm} \quad Remainder: \blank{1.2cm}

    \vspace{2pt}
    The quotient is a difference of squares, so factor completely:
    $x^3+3x^2-4x-12 = $ \blank{5.0cm}

  \item \textbf{Remainder Theorem check.} For that same $f(x) = x^3 + 3x^2 - 4x - 12$, compute
    $f(-3) = $ \blank{1.6cm}

    \vspace{2pt}
    Compare it with the remainder in problem 3. What do you notice? \blank{5.4cm}
\end{enumerate}
\end{tcolorbox}

\begin{tcolorbox}[colback=white, colframe=black!40, title=\textbf{Tier A --- Approaching Proficiency},
    fonttitle=\bfseries, arc=1.5mm, left=3mm, right=3mm, top=2mm, bottom=2mm, breakable]
\begin{enumerate}[leftmargin=1.7em, itemsep=10pt]
  \item \textbf{Long division with a gap.} Divide $x^3 - 5x + 3$ by $x - 2$. Write the dividend with
    its placeholder first.

    \vspace{2pt}
    Dividend in standard form: \blank{6.0cm}

    \vspace{2.0in}

    Quotient: \blank{3.0cm} \quad Remainder: \blank{1.2cm} \quad
    $\dfrac{x^3-5x+3}{x-2} = $ \blank{4.6cm}

  \item \textbf{Synthetic division, then finish the factoring.} Divide $2x^3 + 5x^2 - x - 6$ by
    $x + 2$.
    \begin{center}
    \renewcommand{\arraystretch}{1.6}
    \begin{tabular}{r|C{1.4cm}C{1.4cm}C{1.4cm}C{1.4cm}}
        &     &     &     &     \\
        &     &     &     &     \\ \cline{2-5}
        &     &     &     &     \\
    \end{tabular}
    \end{center}

    \vspace{2pt}
    Depressed polynomial: \blank{3.0cm} \quad Remainder: \blank{1.2cm}

    \vspace{2pt}
    Factor the quotient, then write the complete factorization: \blank{5.4cm}

    \vspace{2pt}
    The three zeros are \blank{3.6cm} \; --- one of them is \emph{not} an integer. Which?
    \blank{1.4cm}

  \item \textbf{No dividing allowed.} Let $g(x) = x^4 - 3x^2 + 2x - 8$. Use the Remainder Theorem
    only.

    \vspace{2pt}
    Remainder when divided by $(x-2)$: \blank{1.6cm} \qquad by $(x+1)$: \blank{1.6cm}

    \vspace{2pt}
    Which of the two is a factor of $g$, and how do you know? \blank{6.0cm}

  \item \textbf{Factor Theorem sweep.} Let $p(x) = x^3 - x^2 - 9x + 9$. Test each candidate.

    \vspace{2pt}
    $p(1) = $ \blank{1.2cm} \quad $p(-1) = $ \blank{1.2cm} \quad $p(3) = $ \blank{1.2cm} \quad
    $p(-3) = $ \blank{1.2cm}

    \vspace{2pt}
    How many factors of the form $(x-r)$ did you just find? \blank{1.2cm} \;
    Write $p$ completely factored: \blank{4.6cm}

    \vspace{2pt}
    This one also factors by \emph{grouping} (Lesson 4.2). Show that route in one line:
    \blank{6.0cm}

  \item \textbf{Error analysis.} Two setups, two different mistakes.

    \vspace{2pt}
    (a) To divide $x^3 - 4x + 1$ by $x - 2$, a student writes the top row as \; $2 \mid 1 \;\; -4
    \;\; 1$. \; What is missing? \blank{4.0cm}

    \vspace{2pt}
    Write the correct top row: \blank{4.6cm} \; and finish: quotient \blank{2.4cm}, remainder
    \blank{1.2cm}

    \vspace{2pt}
    (b) To divide by $x + 5$, a student puts $5$ in the corner box. Why is that wrong, and what
    belongs there? \blank{5.4cm}
\end{enumerate}
\end{tcolorbox}

\begin{tcolorbox}[colback=white, colframe=black!40, title=\textbf{Tier E --- Extension},
    fonttitle=\bfseries, arc=1.5mm, left=3mm, right=3mm, top=2mm, bottom=2mm, breakable]
\textbf{Part 1 --- A divisor synthetic division cannot touch.} Divide
$x^4 + 2x^3 - 7x^2 - 8x + 12$ by $x^2 + x - 6$.
\begin{enumerate}[label=(\alph*), leftmargin=1.8em, itemsep=6pt]
  \item Why can this one \emph{not} be done synthetically? \blank{5.4cm}
  \item Long-divide. \quad Quotient: \blank{3.0cm} \quad Remainder: \blank{1.2cm}

    \vspace{1.6in}

  \item Both the divisor and the quotient factor. Write the original polynomial as a product of four
        linear factors: \blank{6.0cm}
  \item List its four zeros: \blank{4.6cm}
\end{enumerate}

\vspace{5pt}
\textbf{Part 2 --- Solve for the missing coefficient.} Find the value of $k$ that makes $(x-3)$ a
factor of $f(x) = x^3 + kx^2 - 5x + 6$.
\begin{enumerate}[label=(\alph*), leftmargin=1.8em, itemsep=6pt]
  \item By the Factor Theorem, $(x-3)$ is a factor exactly when $f(3) = $ \blank{1.0cm}. Write that
        equation in terms of $k$: \blank{4.6cm}
  \item Solve: $k = $ \blank{1.4cm} \; So $f(x) = $ \blank{4.4cm}
  \item Divide out $(x-3)$ and factor completely: \blank{5.4cm}
\end{enumerate}

\vspace{5pt}
\textbf{Part 3 --- Prove it, then exploit it.}
\begin{enumerate}[label=(\alph*), leftmargin=1.8em, itemsep=6pt]
  \item Start from $f(x) = (x-r)q(x) + R$ and substitute $x = r$ to prove the Remainder Theorem in
        two lines. \writelines{2}
  \item Because the remainder \emph{is} $f(r)$, synthetic division is also a way to \emph{evaluate} a
        polynomial. Use it to find $f(4)$ for $f(x) = x^4 - 3x^3 + 2x - 7$:
    \begin{center}
    \renewcommand{\arraystretch}{1.6}
    \begin{tabular}{r|C{1.3cm}C{1.3cm}C{1.3cm}C{1.3cm}C{1.3cm}}
        &     &     &     &     &     \\
        &     &     &     &     &     \\ \cline{2-6}
        &     &     &     &     &     \\
    \end{tabular}
    \end{center}
    $f(4) = $ \blank{1.6cm} \; Now check by substituting $4$ directly: \blank{4.0cm}
  \item Which method used fewer multiplications? \blank{4.0cm}
\end{enumerate}

\vspace{5pt}
\textbf{Part 4 --- Tomorrow's problem, today.} Nobody will hand you the value of $r$ forever. Take
$h(x) = 2x^3 - 3x^2 - 11x + 6$ and test the integers that divide the constant term $6$.
\begin{enumerate}[label=(\alph*), leftmargin=1.8em, itemsep=6pt]
  \item Candidates: $\pm 1, \pm 2, \pm 3, \pm 6$. Which two give $h(r) = 0$? \blank{3.0cm}
  \item Divide out one of them, say $(x+2)$, and factor the quotient completely:
        $h(x) = $ \blank{5.4cm}
  \item List all three zeros: \blank{3.6cm} \; One of them was \emph{not} on your candidate list.
        Which, and why did the list miss it? \writelines{2}
  \item So what would a complete candidate list have to include? \writeline
\end{enumerate}
\end{tcolorbox}

\end{document}
